% voicetraining.tex
\documentclass{article}

\usepackage{amssymb}
\usepackage{csquotes}
\usepackage{datetime2}
\usepackage{enumitem}
\usepackage{fontspec}
\setmainfont{Atkinson Hyperlegible Next}
\usepackage[margin=1in]{geometry}
\usepackage{multicol}

\pagenumbering{gobble}

\title{Fern's Voice Training!}
\author{}
\date{}

\begin{document}
\maketitle

\section*{Warm-up}
\subsection*{Belly Breathing}
Place one hand on your chest and one hand on your belly.
Breathe in through your nose, you should feel your belly press out and the hand on your chest should stay relatively still.
Then, breathe out with pursed lips.
Continue this breathing method throughout the exercises.

\subsection*{\enquote{Mmm} Sounds}
Hold a hum sound for 5 to 8 seconds.
Try to notice the forward vibration on your lips.
Then, say the following syllables.
\begin{multicols}{5}
    \begin{itemize}
        \item Me Me
        \item May May
        \item My My
        \item Moe Moe
        \item Moo Moo
    \end{itemize}
\end{multicols}

\subsection*{\enquote{Ahh} Sounds}
Take a deep breath and hold a strong, clear \enquote{Ahh} for at least 10 seconds.

\pagebreak
\section*{Phonation}
\subsection*{Easy Onset Practice}
Say the following word pairs.
At the second word, try to mimic the feeling of using air to begin the sound instead of a hard start.
\begin{multicols}{5}
    \begin{center}
        \begin{itemize}
            \item hail, ale
            \item hair, air
            \item hall, all
            \item ham, am
            \item harm, arm
            \item has, as
            \item hear, ear
            \item heat, eat
            \item heel, eel
            \item high, I
            \item hike, Ike
            \item hill, ill
            \item hit, it
            \item hive, I've
            \item how, ow
        \end{itemize}
    \end{center}
\end{multicols}

\subsection*{Easy Onset Sentences}
Say the following sentences with relaxed voicing.
\begin{itemize}
    \item Our Aunt Alice is always angry.
    \item It is impossible to understand arithmetic.
    \item Eryn's umbrella is under the overhang.
    \item Astronauts are unbelievably ambitious.
    \item Everyone asks if Amy is my aunt.
    \item I'll have already eaten enough oatmeal.
    \item Emma is an unusually artistic adolescent.
    \item If you enter early, attendandts are available in the aisles.
    \item Actors arrive at the awards in expensive attire.
    \item Accountants are always extremely accurate.
    \item Erica and Ivana are always at the opera.
    \item Olivia was embarassed by all of Alice's attention.
    \item It was an eye-opening experience for Allison.
    \item It is an effort to understand Aunt Alice's accent.
    \item Olivia's used automobiles are usually under eight thousand dollars.
\end{itemize}

\pagebreak
\section*{Resonance}
As in the Warm-up, start each word with a hum sound and carry it through without letting your voice drop into your throat.
\subsection*{One-syllable Words}
\begin{multicols}{5}
    \begin{itemize}
        \item mad
        \item made
        \item mail
        \item main
        \item make
        \item mall
        \item malt
        \item man
        \item map
        \item mare
        \item mark
        \item mash
        \item mask
        \item match
        \item may
        \item meal
        \item mean
        \item meet
        \item meek
        \item melt
        \item merge
        \item mess
        \item met
        \item mild
        \item mile
        \item milk
        \item mind
        \item mine
        \item mink
        \item mint
        \item mist
        \item mix
        \item mob
        \item mode
        \item mold
        \item mole
        \item mom
        \item monk
        \item month
        \item mood
        \item moon
        \item moose
        \item mop
        \item more
        \item moss
        \item most
        \item moth
        \item mound
        \item mourn
        \item mouse
        \item mouth
        \item move
        \item munch
        \item must
        \item my
    \end{itemize}
\end{multicols}

\subsection*{Two-syllable Words}
\begin{multicols}{4}
    \begin{itemize}
        \item machine
        \item magic
        \item maiden
        \item mailbag
        \item major
        \item mango
        \item manmade
        \item maple
        \item market
        \item marshall
        \item mascot
        \item master
        \item matchbox
        \item measure
        \item meatloaf
        \item melon
        \item melted
        \item member
        \item memo
        \item menace
        \item merger
        \item merit
        \item messy
        \item metric
        \item middle
        \item mighty
        \item million
        \item mindless
        \item minute
        \item misread
        \item mitten
        \item mixer
        \item mobile
        \item model
        \item mohawk
        \item moldy
        \item molten
        \item monarch
        \item money
        \item monkey
        \item monster
        \item moody
        \item moonbeam
        \item morning
        \item mouthpiece
        \item movie
        \item mower
        \item muscle
        \item music
        \item musty
        \item myself
    \end{itemize}
\end{multicols}

\pagebreak
\subsection*{Sentences}
\begin{itemize}
    \item Meet my mom.
    \item Mix mint milk.
    \item Merge more movies.
    \item Moths make a mess.
    \item March for a mile.
    \item Mark my mailbag.
    \item Make maple malt.
    \item Mean marketers mumble.
    \item Munch on mushy melons.
    \item March on Main Street.
    \item Meet the man at the metal mine.
    \item Measure a million miles.
    \item Muddy mowers make a mess.
    \item Master movers never mess up.
    \item Metric machines mean more money.
    \item Mark the message before I move.
    \item Mom melted my magic mascot.
    \item Middle school menaces are mean.
    \item Musicians' mouthpieces make music.
    \item Monsters move in the moonbeams.
    \item My mouse and monkey wear mini mittens.
    \item Members made mobiles with modern matchboxes.
    \item Monarchs love misty mangos in the morning.
    \item Mountain climbers are mystified by monumental peaks.
    \item Mounds of mud messed up the merging motorists.
    \item Moss and mold are not manmade materials.
    \item Markets in Montana collect money from mountaineers.
    \item Mysterious marble is magnified with a microscope.
    \item Marathon runners move quickly on the last mile.
    \item Miniature mice made a home in my mousetrap.
\end{itemize}

\pagebreak
\section*{Intonation}
\subsection*{Statement or Question}
Say each sentence twice, once as a statement and once as a question.
Use falling and rising pitch to differentiate the two respectively.
\begin{multicols}{2}
    \begin{itemize}
        \item You want to go to the park.
        \item I have to pay my bills on Wednesday.
        \item Her mother's name is Susan.
        \item Tomorrow we will have the party.
        \item There is no more room in this cabinet.
        \item It has been six months since we last spoke.
        \item Someone is knocking at the door.
        \item Your prom dress is lavender.
        \item It is too late to call her at home.
        \item Kimberly's cat has had kittens.
    \end{itemize}
\end{multicols}

\subsection*{Words}
For each of the following sentences, read them once in the way that sounds natural.
Then, use pitch inflection to stress each noun as though responding to a question.
\begin{multicols}{2}
    \begin{itemize}
        \item My kitten's name is Casey.
        \item Jane goes to college in Pittsburgh
        \item I have to work at 8 a.m. on Thursday.
        \item Francine's hobby is antique car repair.
        \item I would like a silver bracelet for graduation.
        \item Eryn painted her bedroom blue.
        \item A gold ring will please most any girl.
        \item She saw a cat in the neighbor's house.
        \item Dots of light betrayed the black cat.
        \item She called her name many times.
    \end{itemize}
\end{multicols}

\end{document}
